\section{Background}

\noindent\textbf{Security Monitor:}
\label{sec:securitymonitor}
In many enclave platforms, the security monitor (SM) \cite{costan2016sanctum, lebedev2019sanctorum} is a small (9K LOC in our prototype), trusted piece of software running at a higher privilege mode than the hypervisor or the OS (see Figure \ref{fig:computing_stack}).
Its role is to link low-level invariants exposed by the hardware (e.g., if a specific memory access is authorized), to the high-level security policies defined by the platform (e.g., which enclave can access which memory regions) and to invoke the right microarchitectural cleanup routine.
\\ 

\noindent\textbf{Microarchitectural Timing Side Channels:}
Microarchitectural side channels enable information leakage between security domains by exploiting shared microarchitectural structures. These include memory caches \cite{gotzfried2017cache, brasser2017software}, translation look-aside buffers (TLBs) \cite{gras2018translation}, branch predictors \cite{evtyushkin2018branchscope, aciiccmez2006predicting}, DRAM controllers \cite{wang2014timing}, or any other shared microarchitecture. Attackers typically rely on event timing to indirectly observe microarchitectural state, as it cannot be directly accessed through the instruction set architecture (ISA).
For example, cached data access is faster than main memory access. By manipulating shared microarchitectural state and observing changes due to victim activity, attackers can infer sensitive information about the victim's execution, such as accessed addresses or branch directions.
\\ 

\noindent\textbf{Transient Execution Side Channels (TES):}
\label{sec:specandside}
Speculative or transient execution, a key performance feature in modern processors, complicates security boundary reasoning.
Attackers can manipulate microarchitectural structures that orchestrate speculation (like branch predictors) to trigger speculative execution of code snippets that bypass normal security checks.
Although these erroneous executions are eventually squashed, they alter microarchitectural state (e.g., cache states), potentially leaking sensitive information through side channels.
This class of attacks, known as Spectre, was introduced in Kocher et al. \cite{kocher2018spectre} and has since spawned numerous variants.
\\

\noindent\textbf{Hardware Software Contract for Security:}
Recent work \cite{guarnieri2021hardware, mosier2022axiomatic,cauligi2022sok} has focused on defining formal foundations for speculative security properties using \emph{hardware-software contracts}.
Contract semantics consist of two components: a \emph{leakage model} (\sw{leak}) and an \emph{execution model} (\sw{exec}) and are represented using the notation $\contract[leak]{\cdot}{exec}$.
The leakage model describes what \emph{architectural} state is made public by the software (e.g. execution trace and trace of memory accesses for the \sw{ct} model).
The execution model, on the other hand, models the ``safe'' speculative behavior i.e. which speculative control flows are considered \emph{public} by the software.
A defense mechanism, represented by the \emph{attacker model} $\uarch[]{\cdot}{}$ (formally a \emph{projection} of the microarchitectural state), satisfies a \emph{contract} $\contract[]{\cdot}{}$ if it does not leak more information than what the contract specifies as public.
We write $\uarch[]{\cdot}{} \vdash \contract[]{\cdot}{}$.
Although the software does not directly satisfy contracts per say, it can verify \emph{non-interference} properties with respect to contracts.
For a program $p$, \emph{direct non-interference} with respect to contract $\contract[]{\cdot}{}$ means the program will not make any \emph{secret} information public under that contract.
We write this as $p \vdash NI(\contract[]{\cdot}{})$ and say $p$ satisfies $\contract[]{\cdot}{}$ for simplicity.
Additionally, \emph{relative non-interference} from $\contract[]{\cdot}{A}$ to $\contract[]{\cdot}{B}$, means that it will not expose more information under $\contract[]{\cdot}{B}$ than what is made public by $\contract[]{\cdot}{A}$.
We write $p \vdash NI(\contract[]{\cdot}{A} \Rightarrow \contract[]{\cdot}{B})$.
Detailed definitions are recalled in appendix~\ref{app:hwsfcontract}.